\clearpage
\renewcommand{\sectioncolor}{nano}
\section{How to \emph{find} a proof --- Tao's nine steps}
\markboth{Finding a proof}{}

The techniques in §5 tell you what a finished proof looks like. They do not tell you how to get one.
For that your shelf has Terence Tao's \textit{Solving Mathematical Problems}
(\bk{04\_Problem\_Strategy}), whose Chapter 1 lays out a sequence you can actually follow when stuck.

\begin{center}
\begin{tikzpicture}[font=\scriptsize,
 s/.style={rounded corners=3pt,draw=#1,line width=0.9pt,fill=#1!8,inner sep=6pt,
           text width=4.75cm,align=left,minimum height=1.25cm},
 num/.style={circle,fill=#1,text=white,font=\scriptsize\bfseries,inner sep=2.6pt}]
\foreach \i/\txt/\y/\c in {%
 1/{\textbf{Understand the problem.} What kind is it? Find, prove, or evaluate?}/0/nano,
 2/{\textbf{Understand the data.} What is given, and how do the objects interact?}/-1.45/nano,
 3/{\textbf{Understand the objective.} What exactly must be produced?}/-2.90/nano,
 4/{\textbf{Select good notation.} Represent data and objective as simply as possible.}/-4.35/sand,
 5/{\textbf{Write down what you know}; draw a diagram.}/-5.80/sand}{
 \node[s=\c] (a\i) at (0,\y) {\txt}; \node[num=\c] at (-2.75,\y) {\i};}
\foreach \i/\txt/\y/\c in {%
 6/{\textbf{Modify the problem slightly.} Special cases, relaxed hypotheses.}/0/bio,
 7/{\textbf{Modify the problem significantly.} A more aggressive reformulation.}/-1.45/bio,
 8/{\textbf{Prove results about the question.} Data exists to be used.}/-2.90/bio,
 9/{\textbf{Simplify, exploit data, reach tactical goals.}}/-4.35/bio,
 {}/{\textcolor{proof}{\textbf{Then write it up} using a technique from §5.}}/-5.80/proof}{
 \node[s=\c] at (6.6,\y) {\txt};
 \ifx\i\empty\else\node[num=\c] at (3.85,\y) {\i};\fi}
\end{tikzpicture}
\end{center}
\vspace{2pt}
\begin{center}\begin{minipage}{0.93\textwidth}\footnotesize\color{slate}
\textbf{Figure 2.}\; Tao's sequence, from \textit{Solving Mathematical Problems} ch.~1. Steps 1--5
are setup; 6--9 are attack. \textbf{Read Tao after you have been stuck}, not before --- the advice
only lands once you have felt the need for it.
\end{minipage}\end{center}

\begin{bookbox}[title={Diedrichs \& Lovett's honest warning, §2.3.6}]
\begin{quote}\itshape
``There is \textbf{no universal strategy} to prove or disprove any given mathematical statement. One
of the challenges of advanced mathematics is that the work becomes less and less algorithmic the
further one advances. Proving mathematical statements \textbf{requires creativity}\ldots\ Nonetheless,
there are a few common things to try.''
\end{quote}
Their two suggestions: \textbf{modify an example or other proof}, and \textbf{try a few examples}
(with the faulty-generalization caveat from §3).
\end{bookbox}

\begin{techbox}[title={And the part nobody tells beginners --- D\&L §2.3.6}]
\begin{quote}\itshape
``Though it may take considerable work to find a proof of a statement, \textbf{we do not include in
the proof all of the ideas, mental movements, dead ends, or examples that led up to our proof.}''
\end{quote}
This is why published proofs look effortless and yours will not. \textbf{The finished proof is not a
record of how you found it.} Expect the search to be messy; write the result clean.
\end{techbox}

\begin{dobox}[title={Worked: how §9's theorem was actually found}]
Using Tao's numbering, on the theorem $\operatorname{Tr}(T^N)=\lambda_+^N+\lambda_-^N$:
\begin{center}\small
\begin{tabular}{@{}c>{\raggedright\arraybackslash}p{13.4cm}@{}}
\toprule
\rowcolor{nano!12}
\textbf{step} & \textbf{what happened}\\
\midrule
1--3 & Objective: a formula for $\operatorname{Tr}(T^N)$. Data: $T$ is $2\times2$, real, symmetric.\\
4 & Notation: call the eigenvalues $\lambda_\pm$ rather than carrying $2\cosh\beta$ around.\\
5 & Wrote $T$ out; noticed the symmetry immediately.\\
6 & \textbf{Modified slightly:} tried $N=1,2$ by hand. $\operatorname{Tr}(T)=\lambda_++\lambda_-$ --- suggestive.\\
8 & \textbf{Proved a result about the question:} $T$ symmetric $\Rightarrow$ orthonormal eigenbasis (Axler 7.29). This was the key move.\\
9 & Simplified: in the eigenbasis $T^N$ is trivially diagonal; trace is basis-independent (10.A). Done.\\
\bottomrule
\end{tabular}
\end{center}
\vspace{2pt}
The write-up in §7 is four lines. The search was not.
\end{dobox}
